Let $r$ be the frequency of the unique mode 15. Every other value occurs at most $r-1$ times, and the maximum is at least 28 because the minimum is positive and the range is 27. For an odd sample size $2k+1$, the upper half plus the median sum to at least $13(k+1)+2r+15$. The lower half sum to at least $2k-(r-1)$, so mean eight implies $k\ge r+21$. The alternative lower bound $3k-3(r-1)$ gives $r\ge23$. Thus $k\ge44$ and the sample has at least 89 values. For an even size $2k$, the middle pair must be $(11,15),(12,14)$ or $(13,13)$: its sum is 26 and its upper value cannot exceed the mode 15. The first pair already gives a sum greater than $16k$. For $(12,14)$ the lower and upper sums are at least $2k-r+11$ and $14k+r+14$, again impossible. For $(13,13)$ the bounds $2k-r+12$ and $13k+2r+15$ imply $k\ge r+27$; replacing the first by $3k-3r+13$ gives $r\ge28$, hence size at least 110. Finally the multiset with 22 ones, 22 twos, 21 thirteens, 23 fifteens and one 28 has 89 entries, sum 712, median 13, unique mode 15 and range 27. The minimum is $\boxed{89}$.
